\documentclass[10pt,a4paper]{article}

% ── FONTS ────────────────────────────────────────────────────
\usepackage{fontspec}
\setmainfont{Calibri}
\setsansfont{Calibri}

% ── LAYOUT ───────────────────────────────────────────────────
\usepackage[top=1.0cm,bottom=1.0cm,left=1.5cm,right=1.5cm]{geometry}
\usepackage{graphicx}
\usepackage{xcolor}
\usepackage{booktabs}
\usepackage{tabularx}
\usepackage[inline]{enumitem}
\usepackage{setspace}
\usepackage{fancyhdr}
\usepackage{titlesec}
\usepackage{hyperref}
\usepackage{microtype}
\usepackage{array}
\usepackage{colortbl}
\usepackage{tcolorbox}

\hypersetup{colorlinks=true,linkcolor=navyblue,urlcolor=accentteal,citecolor=navyblue}

% ── BRAND COLORS ────────────────────────────────────────────
\definecolor{navyblue}{HTML}{003566}
\definecolor{gold}{HTML}{C9A84C}
\definecolor{lightgray}{HTML}{F4F6FB}
\definecolor{darkgray}{HTML}{4A4A4A}
\definecolor{accentteal}{HTML}{0077B6}

% ── SECTION STYLING ─────────────────────────────────────────
\titleformat{\section}[block]
  {\normalfont\bfseries\fontsize{10.5}{12.5}\selectfont\color{navyblue}}
  {}{0em}{}
\titlespacing*{\section}{0pt}{6pt}{3pt}

\titleformat{\subsection}[block]
  {\normalfont\bfseries\fontsize{9.5}{11}\selectfont\color{accentteal}}
  {}{0em}{}
\titlespacing*{\subsection}{0pt}{4pt}{1pt}

% ── HEADER / FOOTER ─────────────────────────────────────────
\setlength{\headheight}{15pt}
\pagestyle{fancy}
\fancyhf{}
\fancyhead[L]{\footnotesize\color{navyblue}\textbf{AIGGPA Bhopal} \quad\textbar\quad Forest Department Digitization Assessment}
\fancyhead[R]{\footnotesize\color{navyblue}June 2026}
\fancyfoot[C]{\footnotesize\color{navyblue}\thepage}
\fancyfoot[L]{\scriptsize\color{darkgray}Atal Bihari Vajpayee Institute of Good Governance and Policy Analysis}
\renewcommand{\headrulewidth}{0.4pt}
\renewcommand{\headrule}{\color{navyblue}\hrule width\headwidth height\headrulewidth}
\renewcommand{\footrulewidth}{1pt}
\renewcommand{\footrule}{\color{gold}\hrule width\headwidth height 1pt}

% ── LISTS ───────────────────────────────────────────────────
\setlist[itemize]{leftmargin=1.1em, itemsep=0.5pt, topsep=1pt, parsep=0.5pt}

% ────────────────────────────────────────────────────────────
\begin{document}
\setstretch{1.0}

\begin{center}
  {\fontsize{13}{15}\selectfont\bfseries\color{navyblue}
  Forest Department Survey Data Summary\\[2pt]
  {\fontsize{10.5}{12}\selectfont\bfseries\color{accentteal}Executive Brief \& Strategic Insights on Digital Transformation}}\\[3pt]
  {\small\color{darkgray} AIGGPA Bhopal Summer Research Program \quad\textbar\quad June 2026}
\end{center}

\vspace{-0.2cm}
\noindent
\small The analysis of the Forest Department survey data reveals a workforce generally optimistic about digital transformation, with clear opportunities for further targeted training and system integration.

\vspace{0.1cm}

% ── KEY TAKEAWAYS BOX ───────────────────────────────────────
\begin{tcolorbox}[colback=navyblue!4,colframe=navyblue,leftrule=3.5pt,rightrule=0pt,toprule=0pt,bottomrule=0pt,arc=0pt,top=4pt,bottom=4pt,left=8pt,right=8pt]
\fontsize{9}{12}\selectfont
\noindent{\bfseries\color{navyblue}Key Takeaways}
\begin{itemize}
  \item \textbf{Positive Digital Sentiment:} Respondents strongly believe digital tools make their work faster, improve quality, and increase productivity, with an average utility score of \textbf{3.29 out of 5.0}.
  \item \textbf{Experienced Demographics:} The cohort represents the official department structure with a realistic gender mix (\textbf{75.0\% Male, 25.0\% Female}) and is primarily balanced between mid-career and senior professionals (\textbf{53.8\% aged 30--45, 35.0\% aged 46--60}).
  \item \textbf{High Digital Workload:} A significant proportion of respondents (\textbf{31.2\%}) perform \textbf{61--80\%} of their work digitally, showing high existing integration with official digital systems.
  \item \textbf{Training Gap:} While \textbf{47.5\%} have received formal training, a significant \textbf{52.5\%} have not, representing a critical target area for expanded capacity building programs.
\end{itemize}
\end{tcolorbox}

\vspace{0.05cm}

% ── METRICS SUMMARY TABLE ───────────────────────────────────
\renewcommand{\arraystretch}{1.05}
\begin{table}[h!]
\centering
\footnotesize
\caption{\textbf{Forest Department Survey Cohort \& Digital Usage Profile ($N=80$)}}
\vspace{0.05cm}
\begin{tabularx}{\textwidth}{>{\columncolor{lightgray}}l X >{\columncolor{lightgray}}l X}
\toprule
\rowcolor{navyblue}
\multicolumn{2}{c}{\textcolor{white}{\textbf{Workforce Profile \& Demographics}}} & \multicolumn{2}{c}{\textcolor{white}{\textbf{Digital Capacity \& Tool Adoption}}} \\
\midrule
\textbf{Age Group}       & 30--45: 43 (53.8\%), 46--60: 28 (35.0\%) \par Below 30: 9 (11.2\%) & \textbf{Work digitally} & 61--80\% of workload: 31.2\% \par 81--100\% of workload: 8.8\% \\
\rowcolor{white}
\textbf{Gender Balance}  & Male: 60 (75.0\%) \par Female: 20 (25.0\%) & \textbf{Administrative} & e-Office, Email, MS Office \par Used by 66.3\% (53 out of 80) \\
\textbf{Education Level} & Graduate: 35.0\% \par PG: 22.5\%, 12th: 28.7\%, Prof: 13.8\% & \textbf{Collaboration}  & Google Drive, Zoom, Google Docs \par Used by 47.5\% (38 out of 80) \\
\rowcolor{white}
\textbf{Service Tenure}  & $\ge$ 21 Years: 22.5\% (18), 11--20 Years: 47.5\% (38) \par 6--10 Years: 22.5\% (18), 0--5 Years: 7.5\% (6) & \textbf{Specialized Tools} & e-Green Watch (76.3\%), GIS Portal \par (36.3\%), Forest Alert (45.0\%) \\
\bottomrule
\end{tabularx}
\end{table}

\vspace{-0.25cm}

% ── DETAILED INSIGHTS ───────────────────────────────────────
\section{Workforce Characteristics \& Digital Infrastructure Integration}

\noindent\begin{minipage}[t]{0.49\textwidth}
\subsection{Demographics and Education}
The department's surveyed workforce is highly experienced, well-educated, and structurally balanced:
\begin{itemize}
  \item \textbf{Age:} 43 respondents are aged 30--45 (53.8\%), 28 are aged 46--60 (35.0\%), and 9 are under 30 (11.2\%).
  \item \textbf{Education:} The majority (\textbf{35.0\%}) hold a Graduate degree, with others holding PG degrees (\textbf{22.5\%}), 12th Grade (\textbf{28.7\%}), or Professional degrees (\textbf{13.8\%}).
  \item \textbf{Experience:} A large portion of the staff has significant experience, with \textbf{22.5\%} having over 21 years of service, and \textbf{47.5\%} having 11--20 years of service.
\end{itemize}
\end{minipage}
\hfill
\begin{minipage}[t]{0.49\textwidth}
\subsection{Digital Tool Usage}
The most commonly used digital tools are administrative, with targeted use of domain tools:
\begin{itemize}
  \item \textbf{Administrative Tools:} e-Office, Email, and MS Office are each used by \textbf{66.3\%} of respondents (53 out of 80).
  \item \textbf{Collaboration Platforms:} Google Drive, Zoom, and Google Docs are used by \textbf{47.5\%} of the group (38 out of 80).
  \item \textbf{Specialized Tools:} Specialized platforms like \textbf{e-Green Watch} (76.3\%), \textbf{GIS Portal} (36.3\%), and the \textbf{Forest Alert System} (45.0\%) are noted as vital for field management.
\end{itemize}
\end{minipage}

\vspace{0.05cm}

\section{Challenges \& Strategic Roadmap for Training}
Despite a strong optimistic outlook, respondents highlighted clear obstacles to full digital adoption:
\begin{itemize}
  \item \textbf{Technical \& Infrastructure Issues:} Slow internet connectivity, power outages, and portal security configurations were cited as frequent operational disruptions.
  \item \textbf{Training Deficits:} 42 out of 80 respondents (\textbf{52.5\%}) have not received formal training. Trained respondents rated quality moderately, suggesting that future training sessions should be highly practical and targeted to specific job roles rather than generic.
  \item \textbf{Desired Recommendations:} Open-ended feedback indicates a strong collective preference for a \textbf{"Unified simple administrative portal"} and regular workflow updates to simplify citizen-facing operations.
\end{itemize}

\vspace{0.1cm}

% ── VISUALIZATIONS BOX ──────────────────────────────────────
\begin{tcolorbox}[colback=lightgray!80,colframe=gold,arc=2pt,title=\textbf{Visualizations Roadmap \& Charts Appendix},coltitle=navyblue,fonttitle=\bfseries\footnotesize,top=3pt,bottom=3pt,left=6pt,right=6pt]
\fontsize{8.5}{11}\selectfont
The accompanying analytics dashboard contains the following interactive charts summarizing these findings:
\begin{enumerate*}[label=\textbf{\arabic*.}]
  \item \textbf{Distribution by Age Group}: Shows the staff concentration in the active 30--45 range. \quad
  \item \textbf{Average Digital Sentiment Scores}: Illustrates the high perceived value of digital tools. \quad
  \item \textbf{Training Coverage}: Highlights the gap between trained (47.5\%) and untrained (52.5\%) staff. \quad
  \item \textbf{Digital Workload Distribution}: Displays the extent of existing digital workflow adoption. \quad
  \item \textbf{Commonly Used Digital Tools}: Ranks administrative and collaboration platforms by adoption rate.
\end{enumerate*}
\end{tcolorbox}

\end{document}
